\documentclass[11pt]{article}

\usepackage[margin=1in]{geometry}
\usepackage{amsmath,amssymb,amsthm}
\usepackage{enumitem}
\usepackage{booktabs}
\usepackage{longtable}
\usepackage{xcolor}
\usepackage[hidelinks]{hyperref}
\usepackage{microtype}
\usepackage{listings}

\lstset{
  basicstyle=\ttfamily\small,
  columns=fullflexible,
  breaklines=true,
  frame=single,
  framerule=0.3pt,
  xleftmargin=0.5em,
  xrightmargin=0.5em
}

\newcommand{\ALCIRCC}{\ensuremath{\mathcal{ALCI}_{\mathrm{RCC5}}}}
\newcommand{\Cl}{\operatorname{Cl}}
\newcommand{\Rel}{\operatorname{Rel}}
\newcommand{\EQ}{\mathsf{EQ}}
\newcommand{\PP}{\mathsf{PP}}
\newcommand{\PPI}{\mathsf{PPI}}
\newcommand{\PO}{\mathsf{PO}}
\newcommand{\DR}{\mathsf{DR}}
\newcommand{\blk}{\equiv_{\mathrm{blk}}}
\newcommand{\eqrel}{\equiv_{\EQ}}
\newcommand{\tp}{\operatorname{tp}}
\newcommand{\reach}{\operatorname{reach}}
\newcommand{\dom}{\operatorname{dom}}
\newcommand{\st}{\operatorname{st}}
\newcommand{\PASS}{\textcolor{black!70!green}{\textbf{PASS}}}
\newcommand{\OPEN}{\textcolor{black!60!orange}{\textbf{open}}}

\title{Response to GPT-5.5's Round-2 Review and Verification Recommendations\\
       \large Status report on Python work packages WP1--WP6 and Lean targets L1--L7}
\author{Prepared by Claude (Anthropic), prompted by Michael Wessel}
\date{2026-05-13}

\begin{document}
\maketitle

\begin{abstract}
This document is a response to GPT-5.5's round-2 formal review of the
v2.1 completeness-extraction manuscript and to the subsequent
\emph{verification recommendations} memorandum.  It follows the
expected final-report structure (summary verdict, proof-obligation
table, executable artifacts, Lean artifacts, counterexamples,
remaining gaps, recommended manuscript edits).  The verification work
on the Python side (WP1--WP6) is complete and every script reports
\PASS{}; no counterexample to the repaired proof was found.  The
Lean targets L1--L7 are not yet attempted.  The repaired proof is
adopted in the project as the current completeness pillar.
\end{abstract}

%% ====================================================================
\section{Summary verdict}
%% ====================================================================

\textbf{Conditionally accepted, pending Lean formalization of the
core lemmas.}

The repaired split-forest decidability proof
(\texttt{repaired\_split\_forest\_no\_automata\_proof.tex},
May~2026, 16~pages, no-automata route) has been adopted in the
project as the current completeness pillar, superseding Claude's
earlier completeness-extraction attempts (v1 and v2.1).  The
three load-bearing repairs identified in the round-2 review
\begin{itemize}[leftmargin=2.2em]
  \item \textbf{G1.}  Equality at the level of occurrences,
    \(u \eqrel v \iff \mathrm{orig}(u) = \mathrm{orig}(v)\),
    replacing the v2.1 profile/port equality that contradicted
    vertical separation on self-loops via the M1+M7 reflexivity
    chain.
  \item \textbf{G3.}  Equality mates may have different selected
    parents.  The two-mate stress concept
    \(C_{AB} \equiv A \sqcap B \sqcap \exists\PP.A \sqcap \exists\PP.B
                  \sqcap \forall\PP.(\neg A \vee \neg B)\)
    is accepted with two mate occurrences at distinct parents, and
    rejected when forced onto a single parent chain.
  \item \textbf{G4.}  Rootless orientation:\@ parent self-loops are
    permitted; no ambient root is required.  Request-closed cycles
    replace \(\omega\)-acceptance as the discipline for transitive
    eventualities.
\end{itemize}
all pass the executable tests detailed below.  The remaining
sub-load-bearing items in G2/G5--G8 either no longer apply (because
the underlying construction has changed) or are not contradicted by
any test we ran.

We do \emph{not} recommend further manuscript repairs at this time.
The condition under which the proof is \emph{fully} established is
the Lean formalization of the L1--L7 core lemmas (in particular L1,
L3, L4, L6); we have not attempted these.

%% ====================================================================
\section{Proof-obligation table}
%% ====================================================================

For each of the nine obligations from Section~3 of the verification
recommendations, the table below reports status (proved by GPT-5.5
in the manuscript, tested computationally, counterexample found, or
still open) and the supporting artifact.

\begin{center}\small
\begin{longtable}{p{0.18\linewidth} p{0.40\linewidth} p{0.13\linewidth} p{0.21\linewidth}}
\toprule
Obligation & Statement (abridged) & Status & Evidence \\
\midrule
\endfirsthead
\toprule
Obligation & Statement (abridged) & Status & Evidence \\
\midrule
\endhead
\bottomrule
\endfoot

RCC5 relation algebra
  & Manuscript composition table correct over finite-set semantics;
    \(\PP\circ\PP \subseteq \{\PP\}\), \(\PPI\circ\PPI \subseteq \{\PPI\}\); table
    stabilizes at \(|U|=5\).
  & \PASS{} (computational)
  & WP1: all 25/25 cells match GPT-5.5's table; brute-force
    enumeration over \(|U| \in \{2,3,4,5\}\). \\

Closure / NNF complement
  & \(\Cl(C_0)\) closed under NNF complement; in particular
    \(\overline{\exists R.D} = \forall R.\overline D\).
  & \PASS{} (manuscript)
  & Proved in the repaired manuscript (Def 2.1, Lemma 2.3);
    not separately tested. \\

Witness thinning
  & \(\mathcal I, a_0 \models C_0\)
    \(\Rightarrow\) \(\mathcal I\upharpoonright W, a_0 \models C_0\)
    on the selected-witness substructure.
  & \PASS{} (manuscript)
  & Proved in the repaired manuscript (Lemma 3.2);
    not separately tested. \\

Generated cover (not Hasse)
  & Tree edges are selected generated containment-cover edges; semantic
    \(\PP/\PPI\) is the strict ancestor/descendant closure.
  & \PASS{} (manuscript)
  & Definition 4.1 of the repaired manuscript; no test could
    contradict the choice of orientation. \\

Blocking vs.\ semantic equality
  & \(\blk\) (finite blocking/profile repetition) is kept disjoint
    from \(\eqrel\) (semantic equality / weak-EQ mate relation).
  & \PASS{} (computational)
  & WP2 tests~[1][8] (C\(_{\mathrm{up}}\) self-loop accepted) and
    WP2 test~[3] (G1 reproducer rejected) jointly confirm the
    separation. \\

Multiple upward \(\PP\) witnesses
  & A semantic element may need several incomparable superpart
    witnesses, discharged via explicit \(\eqrel\)-mate
    occurrences with different parents.
  & \PASS{} (computational)
  & WP2 test~[4] accepts \(C_{AB}\) with two mate occurrences;
    WP2 test~[5] rejects the single-parent variant via V4.
    Folds WP4 into WP2. \\

Typed \(\eqrel\) quotient
  & Quotient by typed-RCC5 \(\eqrel\) congruence preserves frame
    conditions and concept truth.
  & \PASS{} (manuscript)
  & Definition 4.3 and Lemma 4.4 of the repaired manuscript; not
    separately tested.  L4 in the Lean stack. \\

Support-closed mosaics
  & Mosaics record joint realizability; every triple in the
    unfolding is covered by a compatible 3-mosaic; larger prefixes
    by overlap-compatible amalgamation.
  & \PASS{} (computational)
  & WP6 finds no counterexample to (M1)--(M5) and (SC1)--(SC4) on
    triples; the (M3) check is exact (41/64 triangles
    consistent, all realizable on \(|U|\le 5\)). \\

Request-closed cycles
  & Cyclic blocking is sound only if every transitive
    \(\PP/\PPI\) existential request is fulfilled later on the
    same finite cycle or via an explicit equality-mate route.
  & \PASS{} (computational)
  & WP2 test~[8] (cycle with \(\exists\PP.A\) discharged by
    self-loop) accepted; WP2 test~[9] (cycle carrying
    \(\exists\PP.A\) but no state satisfies \(A\)) rejected by V1.
    Folds WP5 into WP2. \\

\end{longtable}
\end{center}

\paragraph{Note on coverage.}
\emph{``\PASS{} (manuscript)''} means the obligation is proved in
the repaired manuscript and our verification work did not produce
an additional executable cross-check.  \emph{``\PASS{}
(computational)''} means the corresponding finite-counterexample
search ran and returned no counterexample.  No obligation has
status \emph{counterexample found}.

%% ====================================================================
\section{Executable artifacts}
%% ====================================================================

The Python verification suite is organized as four scripts that
cover all six work packages.  WP4 (multiple-superpart stress) is
folded into WP2 via the dedicated \(C_{AB}\) test; WP5
(request-closed blocking cycles) is folded into WP2 via the cycle
tests, with WP3 providing the small-model cross-check.  The
directory layout follows the suggested structure but with WP4 and
WP5 inlined.

\begin{lstlisting}
verification/
  python/
    rcc5_composition_check.py     WP1
    certificate_checker.py        WP2 + WP4 + WP5  (vertical fragment)
    small_model_oracle.py         WP3
    mosaic_closure_search.py      WP6
  reports/
    rcc5_composition_table.txt
    certificate_checker.txt
    small_model_oracle.txt
    mosaic_closure_search.txt
\end{lstlisting}

\subsection{WP1: RCC5 composition table}

\textbf{Script:}\quad \texttt{rcc5\_composition\_check.py}\\
\textbf{Method:}\quad Compute the RCC5 composition table by
exhaustive enumeration over non-empty subsets of a finite universe
\(U\), using the set-theoretic semantics
\[
A\EQ B \iff A=B,\quad
A\PP B \iff A\subsetneq B,\quad
A\PPI B \iff B\subsetneq A,
\]
\[
A\DR B \iff A\cap B=\varnothing,\quad
A\PO B \iff A\cap B\neq\varnothing\wedge A\nsubseteq B\wedge B\nsubseteq A.
\]
For each \((R,S)\), compute all \(T\) realized by a witness triple
\(A R B\), \(B S C\), \(A T C\).  Increase \(|U|\) until the table
stabilizes.

\textbf{Result:}
\begin{itemize}[leftmargin=2em]
  \item Stability across \(|U| \in \{2,3,4,5\}\):\ \(|U|=5\) and
    \(|U|=4\) give the same signature; the table is therefore
    stable at \(|U|=5\).
  \item Per-cell comparison against the manuscript table
    (lines~99--114 of \texttt{repaired\_split\_forest\_no\_auto\-mata\_proof.tex}):
    \textbf{0 mismatches out of 25 cells}.
  \item Comparison against the in-repo quasimodel reasoner's COMP
    table differs only by the four documented EQ-omission cells
    \((\DR,\DR),(\PO,\PO),(\PP,\PPI),(\PPI,\PP)\)\,---\,a
    convention of that reasoner (which models edges between
    distinct nodes only), not a bug.
\end{itemize}

\textbf{Verdict:}\quad \PASS{}.

\subsection{WP2 + WP4 + WP5: certificate-soundness fuzzer}

\textbf{Script:}\quad \texttt{certificate\_checker.py}\\
\textbf{Coverage:}\quad Validity conditions (V1), (V2), (V3),
(V4), (V5), (V6), (V7), (V8) basic, (V9) Ea/Eb, (V10) of
Definition~4.5 of the repaired manuscript.  The full (V8)
triple-support search and the patchwork closure of (V2)/(V6)/(V7)
are stress-tested separately by WP6.

The fuzzer exercises the three load-bearing repairs:
\begin{itemize}[leftmargin=2em]
  \item \textbf{G1 toggle (occurrence-level equality).}
    The self-loop on
    \(C_{\mathrm{up}} \equiv \exists\PP.\top \sqcap \forall\PP.\exists\PP.\top\)
    is \emph{accepted} as VALID under occurrence-level equality
    (tests~[1][2][8][21][23]) and \emph{rejected} as INVALID when
    two states share an eq-port across a \(\PP\)-edge or across a
    cycle (test~[3], test~[22] V9.Eb).
  \item \textbf{G3 toggle (mates with different parents).}
    The two-mate stress \(C_{AB}\) is \emph{accepted} when the two
    mate occurrences have distinct selected parents (test~[4]),
    and \emph{rejected} (V4) when forced onto a single parent
    chain (test~[5]).  The strengthened variant \(C_{AB}^{\mathrm{inc}}\)
    (cheap-win~1 from the round-2 review) confirms both verdicts
    on a concept with no DAG model (tests~[6][7][8]).
  \item \textbf{G4 toggle (rootless / request-closed).}
    The self-loop discharging \(\exists\PP.A\) at a state carrying
    \(A\) is accepted (test~[11]); the same self-loop without a
    state carrying \(A\) is rejected (test~[12], V1 Hintikka);
    and the pure request-closure variant (cheap-win~2) where
    the self-loop carries \(\neg A \sqcap \exists\PP.A\) and lacks
    an explicit mate is rejected by V4 (test~[13]).
\end{itemize}

The corpus also contains four \emph{occurrence-level distinction}
certs (tests~[21]--[24]) and four \emph{DR/PO mosaic} certs
(tests~[15]--[20]) that exercise V2, V6, V7, V9.Ea, and V9.Eb on
positive and negative inputs.

\textbf{Per-test summary} (corpus of 24 tests):
\begin{lstlisting}
   expected=V  got=V  C_up (PP self-loop, no eq ports)
   expected=V  got=V  C_up with ports, no eq edges
   expected=I  got=I  G1 reproducer: PP plus eq across vertical (V9.Eb)
   expected=V  got=V  C_AB (G3): two mates with different parents
   expected=I  got=I  C_AB v2.1-style single-parent (V4)
   expected=V  got=V  C_AB^inc (G3 strengthened): two mates
   expected=I  got=I  C_AB^inc single-parent attempt (V4)
   expected=I  got=I  C_AB^inc single-chain attempt (V3)
   expected=I  got=I  UNSAT: exists PP.A & forall PP.~A (V1)
   expected=I  got=I  UNSAT: PP-transitive (V3)
   expected=V  got=V  WP5 SAT: cycle, exists PP.A discharged by loop
   expected=I  got=I  WP5 UNSAT: cycle carries exists PP.A, no A
   expected=I  got=I  WP5 pure request-closure: V4 fails, V1 clean
   expected=V  got=V  SAT: forall PP.A & exists PP.A
   expected=V  got=V  DR witness: exists DR.A discharged (V6 SAT)
   expected=V  got=V  PO witness: exists PO.A discharged (V6 SAT)
   expected=I  got=I  DR witness undeclared: no mosaic entry (V6 FAIL)
   expected=I  got=I  DR witness wrong-relation: L=PO not DR (V6 FAIL)
   expected=I  got=I  DR universal violated: forall DR.A but ~A (V2/V7)
   expected=I  got=I  mosaic PP-PP-DR triangle (V2 FAIL)
   expected=V  got=V  Two-state PP cycle, no eq_ports
   expected=I  got=I  Two-state PP cycle with eq across (V9.Eb)
   expected=V  got=V  PP chain (3 states) sharing lambda
   expected=I  got=I  eq across distinct lambdas (V9.Ea)
\end{lstlisting}

\textbf{Verdict:}\quad \PASS{}.

\subsection{WP3: small-model SAT/UNSAT oracle}

\textbf{Script:}\quad \texttt{small\_model\_oracle.py}\\
\textbf{Method:}\quad Brute-force enumeration of complete RCC5
Kripke labelings \(\rho : \dom \times \dom \to \{\EQ,\DR,\PO,\PP,\PPI\}\)
over domain sizes \(n \in \{1,2,3,4\}\) satisfying
\begin{itemize}[leftmargin=2em]
  \item (J1)~reflexivity: \(\rho(x,x) = \EQ\);
  \item (J2)~inverse: \(\rho(y,x) = \mathrm{INV}[\rho(x,y)]\);
  \item (J3)~triple composition:
    \(\rho(x,z) \in \mathrm{comp}(\rho(x,y), \rho(y,z))\).
\end{itemize}
For each labeling and each atomic assignment, the oracle checks
concept satisfaction at the declared root.

\paragraph{Scope and intended role.}  The oracle's domain is
bounded at \(n=4\), so its role in this verification is
\emph{bounded counterexample search}, not evidence for the
infinite-unfolding theorem of Section~5 of the repaired
manuscript.  Concepts such as
\(\exists\PP.\top \sqcap \forall\PP.\exists\PP.\top\) are
\emph{intentionally outside} the oracle's reach:\ they admit only
infinite regular models on the strong-EQ semantics, and asking for
a finite witness at \(n\le 4\) is a category error.  The oracle
contributes to the campaign in two ways only:\ (a) it independently
falsifies any \emph{finite-witnessed} sat claim by exhibiting a
labeling, and (b) it provides a non-trivial cross-check
``WP2~\textsc{Invalid} \(+\) WP3~\textsc{Sat}'' diagnostic conflict
that would flag a bug in the certificate checker.  No instance of
that conflict has been observed.  Evidence for the unfolding
theorem itself is delivered \emph{indirectly} via the certificate
checker (WP2) and the patchwork search (WP6); the oracle does not,
and cannot, certify regular-tree satisfiability on its own.

\textbf{Result on the WP2 cross-check corpus:}
\begin{itemize}[leftmargin=2em]
  \item \(C_{AB}\) is realized at \(n=3\) (consistent with WP2
    accepting its certificate).
  \item Infinite-model-only concepts (e.g.\
    \(C_{\mathrm{up}} \equiv \exists\PP.\top \sqcap \forall\PP.\exists\PP.\top\))
    correctly report \emph{no finite model up to \(n=4\)}; this is
    consistent with SAT in an infinite frame and with WP2
    accepting the certificate (a parent self-loop on the
    rank-\(d\) quotient unfolds to fresh occurrences).  Per the
    scope above, this is the expected outcome, not a failure mode.
  \item Genuine UNSAT concepts (e.g.\
    \(\exists\PP.A \sqcap \forall\PP.\neg A\)) yield \emph{no
    model up to \(n=4\)} \emph{and} WP2 rejects the corresponding
    certificate.
\end{itemize}

No instance of the diagnostic conflict ``WP2~\textsc{Invalid} \(+\)
WP3~\textsc{Sat}'' (which would indicate a bug) was observed.

\textbf{Verdict:}\quad \PASS{}.

\subsection{WP6: mosaic closure / patchwork lemma counterexample search}

\textbf{Script:}\quad \texttt{mosaic\_closure\_search.py}\\
\textbf{Method:}\quad Empirically verify Lemma~4.6 of the repaired
manuscript (patchwork, citing Renz \& Nebel 1999):\ support-closed
mosaics, closed under the local axioms (M1)--(M5) at every
position and under the support-closure operations (SC1)--(SC4) at
every context, yield globally consistent atomic RCC5 networks on
triples.

\textbf{Tests run:}
\begin{itemize}[leftmargin=2em]
  \item \textbf{Test A.}  Partition all \(4^3 = 64\) ordered
    3-position triangles by (M3) consistency:\ 41/64 consistent,
    23/64 rejected.  Every rejected triangle is also unrealizable
    by finite subsets, so (M3) is exact, not over-approximating.
  \item \textbf{Test B.}  For every (M3)-consistent triangle,
    realize it by finite subsets of \(\{0,\dots,n-1\}\) for some
    \(n \le 5\) (Renz--Nebel patchwork base case).  All 41 are
    realizable.
  \item \textbf{Test C.}  DR/PO side-witness insertion: for
    \(\exists R.B\) with \(R \in \{\DR,\PO\}\) and consistent
    universals at \(p\), \(\tau(q)\) can be chosen satisfying
    (M4) at both ends.
  \item \textbf{Test D.}  Universal propagation through PP-chains:
    a broken PP-chain is caught by (M4); a legal one is accepted.
  \item \textbf{Test E.}  Sibling branching (\(p\PPI s_1\),
    \(p\PPI s_2\), \(s_1\PO s_2\), \(s_1\PPI c\)):\ legal
    \(L(c, s_2)\) candidates are \(\{\DR, \PO, \PP\}\), matching
    \(\mathrm{comp}(\PP, \PO)\) exactly.
  \item \textbf{Test F (arity 4, atomic).}  Enumerate all
    \(4^6 = 4096\) atomic 4-networks, partition by (M3) on every
    triple, and realize each (M3)-consistent labeling by 4 nonempty
    finite subsets.  Of the 4096, 916 are (M3)-consistent.  Of these,
    seven require \(n \ge 6\) (the pattern is three pairwise-\(\DR\)
    positions each \(\PO\) a fourth: three external loners plus
    three shared elements in the fourth position, six atoms total).
    Every (M3)-consistent 4-network is realizable for \(n \le 6\).
  \item \textbf{Test G (overlap amalgamation, arity~3 + arity~3).}
    Two triangles \(M_1 = (a,b,c)\) and \(M_2 = (b,c,d)\) sharing
    edge \(\{b,c\}\) with identical labels on the overlap.  For
    every (M3)-consistent pair, search atomic \(L(a,d)\) making
    the joint 4-mosaic (M3)-consistent on all four triples.  All
    427/427 overlap-compatible triangle pairs amalgamate.  This
    is the patchwork lemma at arity~4:\ local agreement implies a
    globally consistent extension.
\end{itemize}

No counterexample to the patchwork lemma was found.  The arity-4
test extends the empirical envelope from 3-mosaics (Tests~A--E) to
4-mosaics in both directions:\ atomic networks (Test~F) and
overlap-glued networks (Test~G).

\textbf{Verdict:}\quad \PASS{}.

\paragraph{Reproducing the reports.}
From the repository root,
\begin{lstlisting}
cd verification/python
python3 rcc5_composition_check.py    > ../reports/rcc5_composition_table.txt
python3 certificate_checker.py       > ../reports/certificate_checker.txt
python3 small_model_oracle.py        > ../reports/small_model_oracle.txt
python3 mosaic_closure_search.py     > ../reports/mosaic_closure_search.txt
\end{lstlisting}
Each script exits with code \texttt{0} on \PASS{} and \texttt{1} on
any mismatch or counterexample.  All four current reports end with
\texttt{VERDICT: PASS}.

%% ====================================================================
\section{Lean artifacts}
%% ====================================================================

\textbf{None.}  The Lean targets L1--L7 from the verification
recommendations have not yet been attempted.  No Lean file exists
in the project for the repaired proof, and consequently no
\texttt{sorry}-free formalization can be reported.  In the Lean
stack the priority order remains as recommended:
\begin{enumerate}[leftmargin=2em]
  \item L1 (RCC5 relation algebra) --- the only target whose
    statement is also covered by an executable check (WP1).
    Recommended first.
  \item L3 (witness thinning).
  \item L4 (typed \(\eqrel\) quotient soundness).
  \item L6 (request-closed lasso correctness).
  \item L2 (closure / NNF complement / Hintikka types),
    L5 (blocking unfolds to fresh), and L7 (mosaic composition
    lifting) after L1/L3/L4/L6.
\end{enumerate}

%% ====================================================================
\section{Counterexamples}
%% ====================================================================

\textbf{None found.}  No script in WP1, WP2, WP3, or WP6 produced
a counterexample to the repaired manuscript's claims on its
tested input families.  In particular:
\begin{itemize}[leftmargin=2em]
  \item WP1 found no entry in the 25-cell composition table at
    which the manuscript differs from the exhaustive enumeration.
  \item WP2 found no certificate whose validity verdict differs
    from the expected outcome on the 24-concept corpus, including
    the three load-bearing G1/G3/G4 reproducers, the two cheap-win
    cases (strengthened \(C_{AB}^{\mathrm{inc}}\) and pure
    request-closure), and the four occurrence-level distinction
    certs.
  \item WP3 found no instance of the diagnostic conflict ``WP2
    INVALID + WP3 SAT''.
  \item WP6 found no triangle for which support-closure rules
    fail to imply RCC5-realizability on \(|U| \le 5\), no DR/PO
    side-witness insertion was blocked, no (M3)-consistent
    4-network was unrealizable for \(n \le 6\), and no
    overlap-compatible triangle pair failed to amalgamate.
\end{itemize}

We emphasize that absence of counterexamples in a bounded search is
not a proof:\ all reported \PASS{} verdicts are evidence, not
certificates.  Pass/fail criterion C8 (effectively checkable
certificate language) is met by the manuscript design but is not
itself executable.

%% ====================================================================
\section{Remaining gaps}
%% ====================================================================

The following lemmas remain informal in our verification work, in
the sense that we have no Lean-checked artifact for them; the
manuscript itself provides hand proofs.

\begin{itemize}[leftmargin=2em]
  \item L1.~RCC5 relation algebra.  Computationally validated by
    WP1; not yet formalized.
  \item L2.~Closure, NNF complement, Hintikka types.  Definitionally
    standard; not yet formalized.
  \item L3.~Witness-generated submodel lemma.  Hand-proved in the
    manuscript (Lemma~3.2);\ not yet formalized.
  \item L4.~Typed-\(\eqrel\) quotient soundness.  Hand-proved in
    the manuscript (Lemma~4.4);\ not yet formalized.
  \item L5.~Blocking unfolds to fresh.  Established by the design
    of the cycle-unfolding map in the manuscript; not yet
    formalized.
  \item L6.~Request-closed lasso correctness.  Hand-proved in the
    manuscript via the cycle-closure conditions;\ tested by WP2
    (tests~[8][9]);\ not yet formalized.
  \item L7.~Mosaic composition lifting.  Tested by WP6;\ not yet
    formalized.  The verification recommendations suggest postponing
    L7 until the computational mosaic rules stabilize, which they
    have.
\end{itemize}

Against the pass/fail criteria (C1)--(C8) of the verification
recommendations:
\begin{center}\small
\begin{tabular}{p{0.05\linewidth} p{0.78\linewidth} p{0.08\linewidth}}
\toprule
ID & Criterion & Status \\
\midrule
C1 & RCC5 composition table executable and independently validated & \PASS{} \\
C2 & Witness-generated submodel lemma formalized
     (or proved in formalization-equivalent detail) & manuscript only \\
C3 & Blocking repetition formally separated from semantic \(\EQ\)
     & \PASS{} (WP2 [3][8]) \\
C4 & Typed-\(\EQ\) quotient soundness proved as a standalone lemma
     & manuscript only \\
C5 & Multiple upward \(\PP\)-witness examples accepted by certificate
     rules and sound after quotienting & \PASS{} (WP2 [4][5]) \\
C6 & Request-closed cycles satisfy all transitive eventualities in
     the unfolding & \PASS{} (WP2 [8][9]) \\
C7 & Mosaics have explicit finite arity bound or finite closure
     construction & \PASS{} (WP6) \\
C8 & Decision procedure enumerates a finite effectively checkable
     certificate language & manuscript only \\
\bottomrule
\end{tabular}
\end{center}

By the criterion of \emph{Section~10} of the verification
recommendations (``If C1--C6 are met but C7 remains informal, the
result should be described as a strong proof strategy with one
remaining finite-interface gap.  If C7 is also closed, then the
proof is much closer to a publication-grade decidability theorem.''),
the work satisfies C1, C3, C5, C6, C7 computationally and
C2, C4, C8 by manuscript hand proof.  We are therefore in the
``much closer to publication-grade'' regime, modulo the Lean
formalization.

%% ====================================================================
\section{Recommended manuscript edits}
%% ====================================================================

\emph{None at this time.}  The verification work did not surface
any text in the manuscript that should be replaced or strengthened
beyond what the repaired proof already states.  Three minor
documentation-level suggestions, in decreasing order of usefulness:

\begin{enumerate}[leftmargin=2em]
  \item \textbf{Composition table line annotation.}  At the
    composition table (lines~99--114), a one-line annotation
    pointing to the WP1 cross-check would let an adversarial reader
    locate the executable check immediately.
  \item \textbf{Folding rationale for WP4/WP5.}  The verification
    recommendations document scopes WP4 and WP5 as separate work
    packages.  In the executed verification suite they are folded
    into WP2 (with WP3 as cross-check), because the
    multiple-superpart stress (\(C_{AB}\)) and the request-closed
    cycle tests are most naturally expressed as additional entries
    in the certificate fuzzer's corpus.  If GPT-5.5 produces a
    revised verification recommendations document, this folding
    could be reflected there.
  \item \textbf{Cycle-close clause for the four EQ-admitting
    composition pairs.}  Independent of the manuscript itself,
    Claude's operational cover-tree tableau adopts a strong-EQ
    cycle-close clause in the role-path compatibility check (Phase
    5 / condition (CT5)) admitting \(w = g\) when
    \((\mathrm{INV}[S], \mathrm{INV}[R]) \in
       \{(\DR,\DR), (\PO,\PO), (\PP,\PPI), (\PPI,\PP)\}\).
    This is the strong-EQ analogue of the same composition pattern
    that produces the four ``conventional EQ omission'' cells in
    the quasimodel reasoner's table.  The repaired manuscript is
    not affected by this clause (the manuscript proof does not use
    a tree-tableau);\ we mention it only because GPT-5.5 may want
    to acknowledge the operational fragment in any future
    revision.
\end{enumerate}

%% ====================================================================
\section{Surrounding alignment work}
%% ====================================================================

For situational awareness, since the round-2 review was issued the
following alignment work has been done in the project:
\begin{itemize}[leftmargin=2em]
  \item \texttt{papers/overview\_ALCIRCC5.tex} (overview paper)
    and \texttt{papers/dl2026\_abstract\_ALCIRCC5.tex} (DL~2026
    abstract) were rewritten to adopt the repaired split-forest
    proof as the current completeness pillar.  The earlier
    completeness-extraction papers
    (\texttt{papers/completeness\_extraction\_ALCIRCC5.tex}~=~v1
    and
    \texttt{papers/completeness\_extraction\_ALCIRCC5\_v2.tex}~=~v2.1)
    are explicitly marked as superseded and kept for historical
    reference.  No claim from those papers is propagated forward.
  \item The v1 split-forest paper
    (\texttt{papers/trees/sibling\_interface\_descriptors\_ALCIRCC5\_completed\_eqsync\_canonical\_needpatched.tex})
    is retained as the source of the soundness direction
    (Thms~1.17--1.19), with a delta paper
    (\texttt{papers/sibling\_interface\_descriptors\_ALCIRCC5\_v2.tex})
    fixing the surrounding Def~1.5, Def~1.7, and Thm~1.20-sketch
    via the corrected generated-cover statements.  The proofs of
    1.17--1.19 carry over verbatim.
  \item The Python quasimodel reasoner
    (\texttt{src/alcircc5\_reasoner.py}) is retained only as a
    cross-validation oracle.  It is known-incomplete on
    cyclic-via-symmetric-role concepts (PO-loop, DR-loop, PP-PPI
    cycle, PPI-PP cycle):\ SAT answers are trustworthy, UNSAT
    answers are sound only for tree-model concepts.
  \item No complexity claim above doubly-exponential
    \(B(C_0) = 2^{2^{p(n)}}\) is asserted for the full logic.  The
    coarse manuscript bound is preserved.  No EXPTIME upper bound is
    claimed.
\end{itemize}

%% ====================================================================
\section*{Bottom line}
%% ====================================================================

The Python WP1--WP6 verification campaign is complete and every
script reports \PASS{}.  No counterexample to the repaired
split-forest decidability proof was found on the tested input
families.  The proof is conditionally accepted as the project's
current completeness pillar.  Closing the remaining gap to a
publication-grade decidability theorem requires the Lean L1--L7
formalization, which we have not attempted; the priority order on
that stack is L1, L3, L4, L6, then L2/L5/L7.

\end{document}
